\subsection{微分与导子}

设 \(\mathbf{B} = \mathbf{A}[(\mathbf{X}_i)_{i\in I}]\) ，则由 III，第 569 页，对每个 \(i\in I\) ，存在唯一的 A-导子 \(\mathbf{D}_i\) （ \(\mathbf{B}\) 的），使得

\[
\mathrm{D} _ {i} \mathrm{X} _ {i} = 1, \quad \mathrm{D} _ {i} \mathrm{X} _ {j} = 0 \quad \text { for } \quad j \neq i.\tag{5}
\]

多项式 \(D_{i}P\) 称为 P 关于 \(X_{i}\) 的偏导数；我们也把它记作 \(D_{X_{i}}P\) 或 \(\frac{\partial P}{\partial X_{i}}\) 或 \(P'_{X_{i}}\) 。由 III，第 558 页，公式 (21)，对于 \(\nu = (\nu_{j}) \in \mathbf{N}^{(I)}\) 我们有

\[
\mathrm{D} _ {i} \left(\mathbf {X} ^ {\nu}\right) = \left\{ \begin{array}{c c c} \nu_ {i} \mathbf {X} _ {i} ^ {\nu_ {i} - 1} \prod_ {j \in I - \{i \}} \mathbf {X} _ {j} ^ {\nu_ {j}} & \text {if} & \nu_ {i} > 0 \\ 0 & \text {if} & \nu_ {i} = 0. \end{array} \right.\tag{6}
\]

由 (6) 可知 \(\mathbf{D}_i\mathbf{D}_j = \mathbf{D}_j\mathbf{D}_i\) （对任意 \(i, j \in I\) ）。对于 \(\nu = (\nu_i)_{i \in I} \in \mathbf{N}^{(I)}\) ，我们令 \(\mathbf{D}^{\nu} = \prod_{i \in I} \mathbf{D}_i^{\nu_i}\) 及 \(\nu! = \prod_{i \in I} (\nu_i!)\) 。在 \(\mathbf{N}^{(I)}\) 上取乘积序，我们有

\[
\mathrm{D} ^ {\nu} \left(\mathrm{X} ^ {\mu}\right) = \left\{ \begin{array}{c l} \frac {\mu !}{(\mu - \nu) !} \mathrm{X} ^ {\mu - \nu} & \text { if } \quad \nu \leqslant \mu , \\ 0 & \text { if   not }. \end{array} \right.
\]

当 P 是单个未定元 X 的多项式时， P 的唯一偏导数记作 DP 或 \(\frac{dP}{dX}\) 或 \(P'\) ，并简称为 P 的导数。

再次令 \(\mathbf{B} = \mathbf{A}\left[(\mathbf{X}_i)_{i\in I}\right]\) ；由 III，第 569 页，B 的 A-微分所构成的 B-模 \(\Omega_{\mathbf{A}}(\mathbf{B})\) 以族 \((d\mathbf{X}_i)_{i\in I}\) （即 \(\mathbf{X}_i\) 的微分所成的族）为基。设 \(\partial_i\) 为指标为 \(i\) 的坐标形式（相对于 \(\Omega_{\mathbf{A}}(\mathbf{B})\) 上的这组基）。于是， B 到自身的映射 \(u\mapsto \langle \partial_i,du\rangle\) 是 B 的一个导子，它把 \(\mathbf{X}_i\) 映为 1 ，把 \(\mathbf{X}_j\) 映为 0 （对于 \(j\neq i\) ），因而就是 \(\mathbf{D}_i\) ；换言之，我们有

\[
d u = \sum_ {i \in I} \left(\mathrm{D} _ {i} u\right) d \mathrm{X} _ {i}\tag{7}
\]

对每个 \(u \in \mathbf{B}\) 。若 \(\mathbf{I}\) 是有限的，则 \((\mathbf{D}_i)_{i \in \mathbf{I}}\) 是由 \(\mathbf{B}\) 的导子组成的 B-模的一组基。

命题4. --- 设 E 是一个结合的、交换的且含单位元的 A-代数， \(\mathbf{x} = (x_i)_{i\in I}\) 是 E 的元素族， \(u\) 是 \(\mathbf{A}[(\mathbf{X}_i)_{i\in I}]\) 的元素，且 \(y = u(\mathbf{x})\) 。那么对于从 E 到某个 E-模中的每个导子 D，我们有

\[
\mathrm{D} y = \sum_ {i \in I} \left(\mathrm{D} _ {i} u\right) (\mathbf {x}) \cdot \mathrm{D} x _ {i}.
\]

只需在 \(u\) 为单项式的情形下证明该命题，而在该情形下，结论由 III，第 558 页，命题 6 可得。

推论。——设 \(f \in \mathbf{A}[\mathbf{X}_1, \ldots, \mathbf{X}_p]\) 与 \(g_i \in \mathbf{A}[\mathbf{Y}_1, \ldots, \mathbf{Y}_q]\) （对于 \(1 \leqslant i \leqslant p\) ），并记 \(h = f(g_1, \ldots, g_p)\) ，则对于 \(1 \leqslant j \leqslant q\) 我们有

\[
\frac {\partial h}{\partial \mathbf {Y} _ {j}} = \sum_ {i = 1} ^ {p} \mathrm{D} _ {i} f (g _ {1}, \dots , g _ {p}) \cdot \frac {\partial g _ {i}}{\partial \mathbf {Y} _ {j}}.\tag{8}
\]

这是命题 4 当 \(\mathbf{E} = \mathbf{A}[\mathbf{Y}_1,\dots ,\mathbf{Y}_q]\) 、 \(x_{i} = g_{i}\) 且 \(\mathbf{D} = \partial /\partial \mathbf{Y}_j\) 时的特殊情形。

设 \(\mathbf{X} = (\mathbf{X}_i)_{i\in I}\) ， \(\mathbf{Y} = (\mathbf{Y}_i)_{i\in I}\) 是两族不相交的未定元，并把 \(\mathbf{X} + \mathbf{Y}\) 记为族 \((\mathbf{X}_i + \mathbf{Y}_i)_{i\in I}\) 。给定 \(u\in \mathbf{A}[\mathbf{X}]\) ，考虑元素 \(u(\mathbf{X} + \mathbf{Y})\) （属于 \(\mathbf{A}[\mathbf{X},\mathbf{Y}]\) ）。对于 \(\nu \in \mathbf{N}^{(I)}\) ，我们用 \(\Delta^{\nu}u\) 表示 \(\mathbf{Y}^{\nu}\) 在 \(u(\mathbf{X} + \mathbf{Y})\) 中的系数（把后者视为关于 \(\mathbf{Y}_i\) 的、系数在 \(\mathbf{A}[\mathbf{X}]\) 中的多项式）。根据定义，我们有 \(\Delta^{\nu}u\in \mathbf{A}[\mathbf{X}]\) ，且

\[
u (\mathbf {X} + \mathbf {Y}) = \sum_ {\nu} \left(\Delta^ {\nu} u\right) (\mathbf {X}) \mathbf {Y} ^ {\nu}.\tag{9}
\]

（此处以及本号其余部分，求和均对指标集 \(\mathbf{N}^{(I)}\) 进行，除非另有说明。）

设 \(\mathbf{a} \in \mathbf{A}^{I}\) ，然后在 (9) 中以 \(\mathbf{a}\) 代替 \(\mathbf{X}\) 、以 \(\mathbf{X} - \mathbf{a}\) 代替 \(\mathbf{Y}\) ，我们得到

\[
u (\mathbf {X}) = \sum_ {\nu} \left(\Delta^ {\nu} u\right) (\mathbf {a}) (\mathbf {X} - \mathbf {a}) ^ {\nu}.\tag{10}
\]

特别地，我们有

\[
u (\mathbf {X}) = \sum_ {\nu} \left(\Delta^ {\nu} u\right) (0) \mathbf {X} ^ {\nu}.\tag{11}
\]

如果 \(u, v \in \mathbf{A}[\mathbf{X}]\) ，我们有

\[
\begin{array}{r l} (u v) (\mathbf {X} + \mathbf {Y}) & = \left(\sum_ {\nu} (\Delta^ {\nu} u) (\mathbf {X}) \mathbf {Y} ^ {\nu}\right) \left(\sum_ {\rho} (\Delta^ {\rho} v) (\mathbf {X}) \mathbf {Y} ^ {\rho}\right) \\ & = \sum_ {\sigma} \left[ \sum_ {\nu + \rho = \sigma} (\Delta^ {\nu} u) (\mathbf {X}) (\Delta^ {\rho} v) (\mathbf {X}) \right] \mathbf {Y} ^ {\sigma} \end{array}
\]

因此

\begin{enumerate}
\def\labelenumi{(\arabic{enumi})}
\setcounter{enumi}{11}
\tightlist
\item
\end{enumerate}

\[
\Delta^ {\sigma} (u v) = \sum_ {\nu + \rho = \sigma} (\Delta^ {\nu} u) (\Delta^ {\rho} v).
\]

设 \(\mathbf{Z} = (Z_i)_{i\in I}\) 是另一族未定元。

我们有：

\[
\begin{array}{r l} \sum_ {\nu} (\Delta^ {\nu} u) (\mathbf {X}) (\mathbf {Y} + \mathbf {Z}) ^ {\nu} & = u (\mathbf {X} + \mathbf {Y} + \mathbf {Z}) \\ & = \sum_ {\sigma} (\Delta^ {\sigma} u) (\mathbf {X} + \mathbf {Y}) \mathbf {Z} ^ {\sigma} \\ & = \sum_ {\rho , \sigma} (\Delta^ {\rho} \Delta^ {\sigma} u) (\mathbf {X}) \mathbf {Y} ^ {\rho} \mathbf {Z} ^ {\sigma}, \end{array}
\]

因此由 I，第 99 页，推论 2：

\[
\Delta^ {\rho} \Delta^ {\sigma} u = \frac {(\rho + \sigma) !}{\rho ! \sigma !} \Delta^ {\rho + \sigma} u.\tag{13}
\]

命题5. --- 对任意 \(u \in \mathbf{A}[\mathbf{X}]\) 及 \(\nu \in \mathbf{N}^{(I)}\) ，我们有

\[
\mathbf {D} ^ {\nu} \boldsymbol {u} = \nu ! \Delta^ {\nu} \boldsymbol {u}.
\]

首先假设 \(\nu\) 的长度为 1 ；则存在 I 的元素 \(i\) ，使得 \(\nu = \varepsilon_i\) ，即 \(\nu_i = 1\) 且 \(\nu_j = 0\) （对 I 中所有 \(j \neq i\) ）。公式 (12) 表明 \(\Delta^{\varepsilon_i}\) 是 A-代数 A{[}X{]} 的一个导子，它显然在 \(\mathbf{X}_j\) （对于 \(j \neq i\) ）上取值为零，而在 \(\mathbf{X}_i\) 上取值为 1 。于是我们有 \(\Delta^{\varepsilon_i} = \mathbf{D}_i\) （对每个 \(i \in I\) ）。

由 (13) 我们有

\[
(\rho ! \Delta^ {\rho}) \cdot (\sigma ! \Delta^ {\sigma}) = (\rho + \sigma)! \Delta^ {\rho + \sigma}\tag{14}
\]

此式在 A-模 \(\mathbf{A}[\mathbf{X}]\) 的自同态代数中成立。现在对 \(\nu\) 的长度作归纳，即得 \(\nu! \Delta^{\nu} = \mathbf{D}^{\nu}\) 。

如果 \(\mathbf{A}\) 是一个 \(\mathbf{Q}\) -代数，公式 (9)、(10)、(11) 就可以写成

\begin{enumerate}
\def\labelenumi{(\arabic{enumi})}
\setcounter{enumi}{14}
\tightlist
\item
\end{enumerate}

\[
u (\mathbf {X} + \mathbf {Y}) = \sum_ {\nu} \frac {1}{\nu !} (\mathrm{D} ^ {\nu} u) (\mathbf {X}) \mathbf {Y} ^ {\nu}
\]

\[
u (\mathbf {X}) = \sum_ {\nu} \frac {1}{\nu !} (D ^ {\nu} u) (\mathbf {a}) (\mathbf {X} - \mathbf {a}) ^ {\nu}\tag{16}
\]

\[
u (\mathbf {X}) = \sum_ {\nu} \frac {1}{\nu !} (D ^ {\nu} u) (0) \mathbf {X} ^ {\nu}.\tag{17}
\]

公式(15)、(16)、(17)三者均被称为“泰勒公式”。

命题6（“欧拉恒等式”）。——设 \(u \in \mathbf{A}[\mathbf{X}]\) 是次数为 \(r\) 的齐次多项式；我们有

\[
\sum_ {i \in I} X _ {i} \cdot D _ {i} u = r u.
\]

设 D 为 A{[}X{]} 到自身的 A-线性映射，使得当 v 为 s 次齐次元时 \(D(v) = sv\) 。我们知道（III，第 554 页，例 6）D 是 A{[}X{]} 的一个导子。因此命题 6 是命题 4（IV，第 6 页）的一个推论。

